%=======================================================================
%  Paper V — Section 3 (Itokawa: the one strong inversion signal) FINAL
%
%  Requires the macros listed at the head of Section 2.
%
%  Values superseded from earlier drafts of this section:
%      scan interval 300-3900        -> 300-4000, widened to 300-8000
%      I = 16.888 %                  -> 16.9 %
%      contrast 1.507                -> 1.51 +/- 0.09
%      Dx_COM 16.86 m                -> 16.9 +/- 2.9 m
%      "identical across meshes"     -> bounded at half a scan step
%      compositional bounds          -> now carry their own uncertainties
%      neck model reported as a fit  -> reported as a non-detection
%
%  All values confirmed against paperV_final.json and paper4_recheck.json.
%  Further corrections in this pass:
%      curvature bounds 130 -> 129, neck 261 -> 259 kg/m3
%      neck slab fraction 12.01 -> 12.05 %
%      dipole test moved to the eleven planes below the grain-density
%      ceiling, to match Section 5: g = -0.244, s = -1.141 +/- 0.018, 21 sigma;
%      eleven-plane block added to the sweep table
%=======================================================================

\section{Itokawa: the one strong inversion signal}
\label{sec:itokawa}

Of the five bodies examined, Itokawa is the only one on which the
two-region inversion produces a dispersion improvement that cannot be
accounted for by mesh resolution alone. This section reports that result
and then tries to break it: by changing the mesh, by moving the dividing
plane, by changing the region shape, by changing the assumed mass, and by
rescaling the shape model. The result survives in the sense that the
recovered contrast is always positive and of the same order; it does not
survive in the sense of being independent of the choices made, and
Section~\ref{sec:constrains} takes up what that means.

\subsection{Result at the detected neck}
\label{sec:itokawa:reference}

The neck detector of Section~\ref{sec:methods:neck} returns
$\xsplit = \SI{0.1608}{\kilo\metre}$ with a prominence of $0.170$,
placing $16.43\%$ of the volume in the head, whose centroid lies at
$x_R = \SI{+0.2197}{\kilo\metre}$. This is the reference solution of the
paper. It is not the plane used by \citet{kanamaru2019}, which lies at
$\SI{0.150}{\kilo\metre}$ and encloses $17.89\%$; the two are kept apart
throughout, and Section~\ref{sec:itokawa:comparison} quotes the second
where a comparison with published values is intended.

Scanning the head density over $4000$~kg\,m$^{-3}$ in steps of
$\SI{25}{\kilo\gram\per\cubic\metre}$ and imposing the mass constraint at
each step gives
%
\begin{equation}
\denshead = 2800 \pm \SI{129}{\kilo\gram\per\cubic\metre},
\qquad
\densrest = 1858 \pm \SI{27}{\kilo\gram\per\cubic\metre},
\qquad
C = \denshead/\densrest = 1.51 \pm 0.09 ,
\label{eq:itokawaref}
\end{equation}
%
with a dispersion improvement of $\improve = 16.9\%$ over the uniform
body. The interval on $\denshead$ is the curvature sensitivity interval
defined in Section~\ref{sec:methods:minimisation}, not a confidence
interval; that on $\densrest$ combines it with the plane-placement
scatter of Table~\ref{tab:itokawasweep}; and that on $C$ follows from the
first two, which move in opposite directions because the mass constraint
ties them, so that their fractional contributions add rather than combine
in quadrature: $129/2800 + 25.4/1858 = 4.6\% + 1.4\% = 6.0\%$, the
$\SI{25.4}{\kilo\gram\per\cubic\metre}$ being the curvature interval
propagated through the mass constraint, $[f/(1-f)]\times129$ at
$f = 0.1643$. The $\pm\SI{27}{\kilo\gram\per\cubic\metre}$ quoted for
$\densrest$ above is larger because it adds, in quadrature, the
$\SI{8}{\kilo\gram\per\cubic\metre}$ by which $\densrest$ moves across the
admissible planes; that term belongs to the interval on $\densrest$ but not
to the budget on $C$, which is evaluated at a fixed plane. A further
term comes from the choice of spin axis: taking the rotational potential about
the recovered principal axis instead of the observed pole moves the contrast
by $3.5\%$ (Section~\ref{sec:methods}). That term is a modelling choice
independent of the other two, so it combines with them in quadrature, giving
$6.9\%$; added linearly, as the most conservative reading would have it, the
total is $9.5\%$. We note
that the agreement of the central value with published contrasts,
discussed in Section~\ref{sec:itokawa:comparison}, is well inside this
interval and should not be quoted to a precision it does not support.

The derived quantities are
%
\begin{equation}
\Dmass = (2.75 \pm 0.48)\times10^{9}~\mathrm{kg}
       = 7.7 \pm 1.3\%~\text{of}~\Mtot ,
\qquad
\dCOM = 16.9 \pm \SI{2.9}{\metre} .
\label{eq:itokawaderived}
\end{equation}
%
Both follow from $\densrest$ through
Equations~\eqref{eq:excessidentity} and~\eqref{eq:excessmass} and are not
independent measurements; they are reported because
Section~\ref{sec:itokawa:comparison} compares the second against an
external constraint, and because Section~\ref{sec:constrains} uses the
first. Computed directly from the shape model, the centre of figure lies at
$x_{\mathrm{cf}} = \SI{0.36}{\metre}$, within half a metre of the origin of
the supplied model; inverting Equation~\eqref{eq:excessmass} at this plane
gives $x_R - x_{\mathrm{cf}} = \SI{219.3}{\metre}$ against a region centroid
at $\SI{219.7}{\metre}$, consistent to the rounding.

Two figures should be recorded alongside these. The amplification factor
of Equation~\eqref{eq:amplification} is $\ampfac = 12.0$ at this plane,
so a fractional scatter quoted for $\densrest$ corresponds to twelve
times that in $\Dmass$. And the displacement statistic of
Equation~\eqref{eq:sstat} is $\Sstat = 5.8$, which is the quantity that
separates a detection from a uniform body: for a uniform body
$\densrest \equiv \densbulk$ holds identically, $\Sstat$ is zero however
small the scatter, and no amount of stability in $\densrest$ is evidence
of anything.

\begin{figure}[htbp]
\centering
\includegraphics[width=0.72\textwidth]{fig03_dispersion_curves.pdf}
\caption{Normalized geopotential dispersion against region density, each
body at its own reference plane, normalized by its own bulk density and by
its own uniform-density dispersion so that the five are comparable. Each
body has its own line style and its own marker at the minimum, so the five
remain distinguishable in monochrome. Itokawa has a deep, broad well,
reaching $0.831$ of
the uniform value at a region density $1.39$ times the bulk; the other four
have shallow wells pinned near it, the deepest of them Arrokoth at $0.987$.
Every curve is the run its section reports, the Arrokoth one with the buried
facets of Section~\ref{sec:methods:shapes} excluded, so the improvements in
the legend are those of Table~\ref{tab:nulls}.}
\label{fig:curves}
\end{figure}

\subsection{Resolution stability of the inferred density}
\label{sec:itokawa:mesh}

Paper~IV established that the objective function is stable under
decimation. That does not establish that the density recovered by
minimizing it is stable, which is a separate question and the one that
matters here. We repeated the inversion on meshes of $20{,}000$,
$30{,}000$, $40{,}000$ and $49{,}152$ facets (Figure~\ref{fig:meshconv}).

\begin{figure}[htbp]
\centering
\includegraphics[width=0.92\textwidth]{fig04_mesh_convergence.pdf}
\caption{Recovered density against facet count with the dividing plane held
fixed. Top: Itokawa. The three finest meshes share the plane
$\xsplit = \SI{0.1608}{\kilo\metre}$ and are flat to the resolution of the
scan; the coarsest, which the detector placed at
$\SI{0.1702}{\kilo\metre}$, is shown as a single point and differs by the
plane and not by the resolution. Bottom: Arrokoth at its neck, unchanged
across all three meshes. The shaded band is $\pm$ one scan step, which is
the resolution of the answer; a change smaller than the band is not a
change.}
\label{fig:meshconv}
\end{figure}

The three finest meshes return the same dividing plane,
$\SI{0.1608}{\kilo\metre}$, and on those three the recovered $\denshead$
is unchanged. We state what that does and does not mean. Because no
sub-grid refinement is applied, ``unchanged'' bounds the variation at
half a scan step, $\SI{12.5}{\kilo\gram\per\cubic\metre}$ or $0.45\%$;
it does not resolve it. The accompanying figure, that $\densrest$ moves
by $\SI{0.1}{\kilo\gram\per\cubic\metre}$ or $0.005\%$, is a weaker
statement still: with $\denshead$ pinned to the same grid point and
$\Mtot$ fixed, Equation~\eqref{eq:massbalance} makes $\densrest$ a
function of $\VR$ and $\Vtot$ alone, so the figure measures how little
the enclosed volumes change under decimation rather than how little the
inversion changes. Both are worth reporting; neither is a measurement of
the objective's resolution sensitivity.

The coarsest mesh is the informative case. At $20{,}000$ facets the
detector returns $\SI{0.1702}{\kilo\metre}$ instead of
$\SI{0.1608}{\kilo\metre}$, a displacement of $\SI{9.4}{\metre}$.
Table~\ref{tab:itokawasweep} gives
$\mathrm{d}\denshead/\mathrm{d}\xsplit \simeq
\SI{10.9}{\kilo\gram\per\cubic\metre\per\metre}$ over that interval, so
the displacement alone would move the recovered head density by about
$\SI{100}{\kilo\gram\per\cubic\metre}$, or $3.6\%$---eight times the
bound obtained on the three finer meshes. The value actually returned on
that mesh is $\SI{2900}{\kilo\gram\per\cubic\metre}$, at the plane
$\xsplit = \SI{0.1702}{\kilo\metre}$ that the detector returns there. That is
exactly the value the unreduced mesh returns at
$\SI{0.1700}{\kilo\metre}$ (Table~\ref{tab:itokawasweep}), so the
disagreement between the two meshes is attributable to the plane in full and
to the objective function not at all: each returns the same density at its
own plane, and they differ only because the detector places the plane
$\SI{9.4}{\metre}$ apart on them.
Resolution sensitivity therefore enters this calculation through the
plane detector rather than through the objective function, and a study
that fixed the plane externally would not see it at all. We regard this
as the more useful statement of the two, and it is the reason
Section~\ref{sec:guidelines} recommends reporting the detected plane on
every mesh rather than the density alone.

At the reference plane the improvement of $16.9\%$ is $12.5$ times the
value the envelope of Paper~IV takes on Itokawa specifically, $1.35\%$ over
the ladder from $2{,}000$ to $42{,}000$ facets with quadric decimation. It is
$2.9$ times the upper edge of the cross-body range, $5.9\%$, but that edge is
attained on Eros, so the ratio to it is the weaker of the two comparisons.
The envelope is a scale, not a threshold
(Section~\ref{sec:intro:relaxation}), and the comparison establishes only
that the improvement is not of the size that a change of discretization
would produce.

\subsection{Sensitivity to the dividing plane}
\label{sec:itokawa:sweep}

The plane is chosen, not measured. We therefore swept it from
$\SI{0.130}{\kilo\metre}$ to $\SI{0.240}{\kilo\metre}$ and repeated the
inversion at each position. Table~\ref{tab:itokawasweep} gives the
result. It is the central piece of evidence in this paper, and
Section~\ref{sec:constrains} is built on it.

A statement about the scan is needed before the outer planes are read.
The interval searched was $4000$~kg\,m$^{-3}$ for
$\xsplit \le \SI{0.1800}{\kilo\metre}$ and $8000$~kg\,m$^{-3}$ for
$\xsplit \ge \SI{0.1850}{\kilo\metre}$, the widening being made in
anticipation of the steep rise in recovered density at the outer planes
rather than in response to a boundary minimum: at
$\SI{0.1850}{\kilo\metre}$ the recovered value of $3075$ lies at $38\%$
of its scan limit, whereas at $\SI{0.1800}{\kilo\metre}$, where the
interval was not widened, $3025$ lies at $76\%$ of its own. The two
intervals share the same $\SI{25}{\kilo\gram\per\cubic\metre}$ grid, so
widening adds candidate points without displacing existing ones and no
discontinuity is introduced at the changeover. Every plane in
Table~\ref{tab:itokawasweep} returns an interior minimum under the rule
of Section~\ref{sec:methods:minimisation}; the closest approach to any
scan boundary is $3025$ against a limit of $4000$, a margin of
$\SI{975}{\kilo\gram\per\cubic\metre}$. No row is a boundary minimum and
none is discarded.

Three features of the table should be read before its statistics. First,
the recovered head density rises monotonically and at an accelerating
rate: the gradient $\mathrm{d}\denshead/\mathrm{d}\xsplit$ is
$\SI{7.5}{\kilo\gram\per\cubic\metre\per\metre}$ over the first four
planes and reaches
$\SI{55.0}{\kilo\gram\per\cubic\metre\per\metre}$ over the last two,
while $\densrest$ changes by only
$-\SI{0.22}{\kilo\gram\per\cubic\metre\per\metre}$ over the admissible
range---smaller by a factor of $34$ than the gentlest part of the
$\denshead$ gradient and $250$ than the steepest. Second, the improvement
also rises monotonically, from $12.9\%$ to $26.4\%$, with no interior
turning point: the objective prefers the planes furthest along the body,
which are the planes the composition forbids
(Section~\ref{sec:itokawa:ceiling}). The reported $\improve = 16.9\%$ is
therefore the improvement \emph{at the detected neck}, not the best the
objective can do, and we describe it that way throughout. Interpolating
between the planes at $0.190$ and $\SI{0.195}{\kilo\metre}$, the improvement
at the block-density bound is $21.5\%$: the objective prefers the edge of the
admissible range to the neck by $4.6$ percentage points, a difference that
itself lies inside the mesh-resolution sensitivity envelope. Third,
$\densrest$ is not flat but traces a shallow minimum near
$\SI{0.200}{\kilo\metre}$, falling from $1865.9$ to $1849.8$ and rising
again to $\SI{1855.5}{\kilo\gram\per\cubic\metre}$; the variation is
systematic rather than noise, and the minimum lies outside the admissible
range, so it should not be read as an optimum.

\begin{table}[htbp]
\centering
\small
\setlength{\tabcolsep}{4pt}
\caption{The Itokawa split-plane sweep, thirteen planes on the unreduced
$49{,}152$-facet mesh at $\Mtot = 3.58\times10^{10}$~kg. The planes at
$\SI{0.1600}{\kilo\metre}$ and $\SI{0.1608}{\kilo\metre}$ are distinct,
separated by $\SI{0.8}{\metre}$: the first belongs to the sweep grid, the
second is the saddle returned by the detector and is the reference
solution (bold). The scan column gives the interval over which $\densR$
was searched; it was widened at $\xsplit \ge \SI{0.1850}{\kilo\metre}$ in
anticipation of the steep rise in recovered density, not in response to a
boundary minimum, and the two intervals share the same
$\SI{25}{\kilo\gram\per\cubic\metre}$ grid. Every plane returns an
interior minimum, the closest approach to a boundary being $3025$ against
a limit of $4000$. Planes below the first rule exceed the empirical
block-density bound and those below the second the LL-chondrite
grain-density ceiling of Section~\ref{sec:itokawa:ceiling}; they are
listed because the monotone rise of $\improve$ through them is the
evidence for Section~\ref{sec:plane}, not because they describe a
possible interior. The four quantity columns are ordered by how nearly
invariant they are, and the ordering is the result. $\denshead$ is
recovered on the scan grid with no sub-grid refinement, and the two
places where $\dCOM$ decreases as the plane advances are that grid rather
than a geometric effect, since $x_R - x_{\mathrm{cf}}$ increases
monotonically from $203.9$ to $\SI{262.6}{\metre}$ throughout.}

\label{tab:itokawasweep}
\setlength{\tabcolsep}{10pt}
\begin{tabular}{lrrrrrrr}
\toprule
$\xsplit$ & $\VR/\Vtot$ & $\densrest$ & $\Dmass$ & $\dCOM$
 & $\denshead$ & scan & $\improve$ \\
(km) & (\%) & (kg\,m$^{-3}$) & ($10^{9}$~kg) & (m)
 & (kg\,m$^{-3}$) & (kg\,m$^{-3}$) & (\%) \\
\midrule
$0.1300$ & $20.73$ & $1865.9$ & $2.614$ & $14.89$ & $2575$ & $300$--$4000$ & $12.9$ \\
$0.1400$ & $19.28$ & $1860.7$ & $2.707$ & $15.81$ & $2650$ & $300$--$4000$ & $14.1$ \\
$0.1500$ & $17.89$ & $1857.7$ & $2.761$ & $16.51$ & $2725$ & $300$--$4000$ & $15.4$ \\
$0.1600$ & $16.54$ & $1856.9$ & $2.774$ & $16.97$ & $2800$ & $300$--$4000$ & $16.8$ \\
$\mathbf{0.1608}$ & $\mathbf{16.43}$ & $\mathbf{1858.1}$ & $\mathbf{2.753}$
 & $\mathbf{16.86}$ & $\mathbf{2800}$ & $\mathbf{300}$--$\mathbf{4000}$ & $\mathbf{16.9}$ \\
$0.1700$ & $15.18$ & $1854.1$ & $2.825$ & $17.66$ & $2900$ & $300$--$4000$ & $18.2$ \\
$0.1800$ & $13.80$ & $1850.9$ & $2.881$ & $18.40$ & $3025$ & $300$--$4000$ & $19.6$ \\
$0.1850$ & $13.08$ & $1853.0$ & $2.844$ & $18.37$ & $3075$ & $300$--$8000$ & $20.3$ \\
$0.1900$ & $12.36$ & $1852.5$ & $2.852$ & $18.63$ & $3150$ & $300$--$8000$ & $21.0$ \\
\cmidrule(l){1-8}
\multicolumn{8}{l}{\footnotesize\itshape
 below: $\denshead$ exceeds the empirical block-density bound,
 $3220 \pm \SI{220}{\kilo\gram\per\cubic\metre}$} \\
$0.1950$ & $11.62$ & $1850.2$ & $2.893$ & $19.11$ & $3250$ & $300$--$8000$ & $21.7$ \\
$0.2000$ & $10.87$ & $1849.8$ & $2.901$ & $19.38$ & $3350$ & $300$--$8000$ & $22.4$ \\
\cmidrule(l){1-8}
\multicolumn{8}{l}{\footnotesize\itshape
 below: $\denshead$ exceeds the LL-chondrite grain density,
 $3540 \pm \SI{130}{\kilo\gram\per\cubic\metre}$} \\
$0.2200$ & $7.86$ & $1851.9$ & $2.863$ & $20.04$ & $3900$ & $300$--$8000$ & $24.8$ \\
$0.2400$ & $5.00$ & $1855.5$ & $2.798$ & $20.52$ & $5000$ & $300$--$8000$ & $26.4$ \\
\midrule
\multicolumn{8}{l}{\bfseries All thirteen planes\quad
 ($\VR$ spanning a factor $4.15$)} \\
mean       & $13.90$ & $1855.2$ & $2.805$ & $17.93$ & $3169.2$ & & \\
half-range &         & $0.434\%$ & $5.12\%$ & $15.70\%$ & $38.26\%$ & & \\
CV         &         & $\mathbf{0.251\%}$ & $\mathbf{2.95\%}$
           & $\mathbf{9.31\%}$ & $20.61\%$ & & \\
$\ampfac$  &         & \multicolumn{6}{l}{mean $11.77$, range $11.34$--$12.70$} \\
\midrule
\multicolumn{8}{l}{\bfseries Below the grain-density ceiling, $n=11$\quad 

 ($\VR$ spanning a factor $1.907$) %; used for scaling in Section~\ref{sec:constrains})
 } \\
mean       & $15.25$ & $1855.4$ & $2.800$ & $17.51$ & $2936$ & & \\
half-range &         & $0.434\%$ & $5.12\%$ & $12.82\%$ & $13.20\%$ & & \\
CV         &         & $\mathbf{0.268\%}$ & $\mathbf{3.16\%}$
           & $\mathbf{8.15\%}$ & $8.60\%$ & & \\
$\ampfac$  &         & \multicolumn{6}{l}{mean $11.80$} \\
\midrule
\multicolumn{8}{l}{\bfseries Admissible planes only, $n=9$\quad
 ($\VR$ spanning a factor $1.68$)} \\
mean       & $16.14$ & $1856.6$ & $2.779$ & $17.12$ & $2855.6$ & & \\
half-range &         & $0.404\%$ & $4.80\%$ & $10.92\%$ & $10.07\%$ & & \\
CV         &         & $\mathbf{0.252\%}$ & $\mathbf{3.00\%}$
           & $\mathbf{7.42\%}$ & $6.89\%$ & & \\
$\ampfac$  &         & \multicolumn{6}{l}{mean $11.90$} \\
\bottomrule
\end{tabular}


\vspace{4pt}
\begin{minipage}{\textwidth}
\footnotesize
\textbf{Mass-balance audit.} Every row satisfies
$\Dmass = (\densbulk-\densrest)\Vtot$ with
$\densbulk = \SI{2012.87}{\kilo\gram\per\cubic\metre}$, and independently
$\Dmass = \Ddens\,\VR$, to the printed precision
(Section~\ref{sec:methods:verify}). The $\dCOM$ column is likewise not
independent: by Equation~\eqref{eq:excessmass} it is $\Dmass$ multiplied
by the purely geometric factor $(x_R-x_{\mathrm{cf}})/\Mtot$. Its value
lies in testing a rival hypothesis, not in adding data.
\\[3pt]
\textbf{Exclusion of the dipole hypothesis.} If the conserved quantity
were the mass dipole $D = \Dmass\,(x_R-x_{\mathrm{cf}})$ rather than the
excess mass, $\dCOM$ would be the flatter column. It is not: its
coefficient of variation exceeds that of $\Dmass$ by $2.6$ over the
eleven planes below the grain-density ceiling and $3.2$ over all thirteen.
Quantitatively, $\mathrm{d}\log(x_R-x_{\mathrm{cf}})/\mathrm{d}\log\VR =
-0.244$ over the eleven planes, so dipole conservation requires a scaling
slope $s = -0.756$; the measured value is $s = -1.141 \pm 0.018$,
twenty-one standard errors away (Section~\ref{sec:constrains:slope}).
\end{minipage}
\end{table}

\subsection{A physical bound on the plane}
\label{sec:itokawa:ceiling}

Because the objective rises monotonically with plane position, it cannot
select the plane, and an external bound is needed. Two are available, and
both are bounds on the recovered head density rather than on the plane
directly; each is converted into a plane position through the measured
gradient of Table~\ref{tab:itokawasweep}, and each therefore inherits an
uncertainty that we state rather than suppress.

The hard bound is compositional. Itokawa is an LL~chondrite parent body
\citep{nakamura2011}, and no assemblage of LL material can exceed its
grain density, $3540 \pm \SI{130}{\kilo\gram\per\cubic\metre}$
\citep{consolmagno2008}. A head denser than this would require a
composition the returned samples exclude. Interpolating
Table~\ref{tab:itokawasweep} between $\SI{0.2000}{\kilo\metre}$ and
$\SI{0.2200}{\kilo\metre}$, where the local gradient is
$\SI{27.5}{\kilo\gram\per\cubic\metre\per\metre}$, places the crossing at
$\xsplit = \SI{0.2069}{\kilo\metre}$, and the $\pm130$ on the grain
density makes that $\SI{0.2069+-0.0047}{\kilo\metre}$. Two caveats
attach. The interpolation spans an interval over which
$\denshead(\xsplit)$ is convex---the gradient rises from $20$ to
$\SI{27.5}{\kilo\gram\per\cubic\metre\per\metre}$ across it---so the
secant lies above the curve and the true crossing is at a slightly larger
$\xsplit$: the figure is a conservative lower bound on the admissible
range, not a best estimate. And the grain density is a zero-porosity
limit, so a head at exactly that value would be a solid monolith, which
the surface morphology does not support.

The softer bound is empirical. The bulk density of LL~chondrite falls,
which includes their internal microporosity, is
$3220 \pm \SI{220}{\kilo\gram\per\cubic\metre}$ \citep{consolmagno2008}.
This is a meteorite measurement, not a measurement of boulders on Itokawa,
and it bounds the density of the blocks a rubble pile is built from only on
the assumption that those blocks resemble the falls.
A head region denser than the densest individual block it could be made
of would require the head to be less porous than its own constituents.
The crossing falls between $\SI{0.1900}{\kilo\metre}$ and
$\SI{0.1950}{\kilo\metre}$, where the local gradient is
$\SI{20}{\kilo\gram\per\cubic\metre\per\metre}$, at
$\xsplit = \SI{0.1935}{\kilo\metre}$; the $\pm220$ on the bound makes
that $\SI{194+-11}{\metre}$, which should not be quoted to the tenth of a
metre.

We adopt the softer bound as the edge of the admissible range, giving the
nine planes from $\SI{0.1300}{\kilo\metre}$ to
$\SI{0.1900}{\kilo\metre}$ over which Table~\ref{tab:itokawasweep}
reports its second set of statistics. The detected neck at
$\SI{160.8}{\metre}$ sits comfortably inside it, $\SI{33}{\metre}$ from
the edge and three times the uncertainty on the edge. The essential point
is not the numerical value of either bound but that a bound from outside
the method was needed at all: the objective function, left to itself,
runs to a head density that the meteorite collection forbids.

\subsection{Comparison with published determinations}
\label{sec:itokawa:comparison}

Table~\ref{tab:itokawacompare} sets our result beside the published ones.
Two cautions govern how it should be read.

First, the comparison must be made at a common plane.
\citet{kanamaru2019} divide at $\SI{0.150}{\kilo\metre}$; at that plane
our inversion returns $\denshead = 2725$, $\densrest = 1857.7$,
$C = 1.467$ and $\dCOM = \SI{16.51}{\metre}$, and those are the values
entered in the table. The reference solution of
Section~\ref{sec:itokawa:reference} belongs to a different plane and a
different row, and the two sets should not be interchanged.

Second, three of the four entries are not independent of one another.
\citet{motooka2012}, \citet{kanamaru2019} and this work all minimize a
functional of the surface geopotential, differing in the weighting and in
the surface subset. Agreement among them tests implementations, not
physics, and the recovery of $1.467$ against $1.460$ is expected rather
than confirmatory---the more so since our own interval on $C$ is
$\pm0.09$, two orders of magnitude wider than the difference. The only
entry derived from independent physics is \citet{lowry2014}, whose YORP
spin-up measurement gives $\densbody = 1750 \pm 110$ and
$\denshead = 2850 \pm \SI{500}{\kilo\gram\per\cubic\metre}$, hence a
contrast of $1.63 \pm 0.30$ if the quoted uncertainties are combined in
quadrature, and a centre-of-mass offset of $\SI{21+-12}{\metre}$. Our
$C = 1.47 \pm 0.09$ and $\dCOM = \SI{16.5+-2.9}{\metre}$ lie inside those
intervals, but the intervals are wide, and the test is correspondingly
weak: a contrast anywhere between $1.33$ and $1.93$ would have passed it.

\begin{table}[htbp]
\centering
\footnotesize
\setlength{\tabcolsep}{2pt}
\caption{Published two-density determinations for Itokawa and this work,
all quoted at or referred to the dividing plane
$\xsplit = \SI{0.150}{\kilo\metre}$. The first three rows minimize a
functional of the surface geopotential and are not independent of one
another; only \citet{lowry2014} rests on different physics. Contrasts are
computed from the tabulated densities.}
\label{tab:itokawacompare}
\begin{tabular}{lllrrrr}
\toprule
Source & Method & Surface used & $\denshead$ & $\densbody$ & $C$ & $\dCOM$ \\
       &        &              & \multicolumn{2}{c}{(kg\,m$^{-3}$)} & & (m) \\
\midrule
\citet{lowry2014}       & YORP spin-up            & ---            & $2850\pm500$ & $1750\pm110$ & $1.63\pm0.30$ & $21\pm12$ \\
\citet{kanamaru2019}    & potential variance      & global         & $2730$       & $1870$       & $1.460$       & $16$ \\
\citet{kanamaru2019pss} & potential variance      & smooth terrain & $2450$       & $1930$       & $1.269$       & --- \\
This work               & dispersion minimization & global         & $2725$       & $1857.7$     & $1.47\pm0.09$ & $16.5\pm2.9$ \\
\bottomrule
\end{tabular}
\end{table}
% The centre-of-mass offset of Kanamaru et al. (2019, PSS) is not quoted in
% the same form as the others and is left empty rather than converted.

The spread among the three equipotential-family results is itself
informative. \citet{kanamaru2019pss}, restricting the average to the
smooth terrains, obtain a contrast of $1.269$ against $1.460$ on the
global surface: a change of surface subset moves the answer by more than
our quoted interval. The choice of which facets enter
Equation~\eqref{eq:dispersion} is therefore a modelling decision of the
same order as the choice of dividing plane, and we have not explored it.

\subsection{The head--neck ambiguity}
\label{sec:itokawa:ambiguity}

The head model is one hypothesis about where the anomaly sits. The neck
model of Section~\ref{sec:methods:identity}, in which the dense region is
a slab spanning the constriction, is another, and
\citet{kanamaru2019pss} report that on the smooth terrains it fits
marginally better. On the global surface we find the opposite, and by a
margin that can be stated quantitatively.

The neck slab, with upper edge at $\SI{0.2029}{\kilo\metre}$ and lower
edge at $\SI{0.1187}{\kilo\metre}$, a width of
$w = \SI{0.0842}{\kilo\metre}$, encloses $12.05\%$ of the volume.
The inversion returns
$\densneck = 2525 \pm \SI{259}{\kilo\gram\per\cubic\metre}$ and
$\densrest = \SI{1943}{\kilo\gram\per\cubic\metre}$, a contrast of
$1.300$, with an improvement of $\improve = 2.3\%$. Four comparisons
follow.

The improvement does not clear the bar the paper sets for itself. At
$2.3\%$ it lies inside the $0.9$--$5.9\%$ mesh-resolution envelope, which
is where the four weak-signal bodies of Section~\ref{sec:nulls} also sit.
By the criterion applied to them, the neck model is not a detection, and
we do not report it as one.

The curvature interval is wider in relative terms, $\pm259$ on $2525$
being $10.3\%$ against $4.6\%$ for the head model, so the objective is
less sharply curved about its minimum and the parameter is less well
determined.

The displacement statistic is smaller and its denominator larger.
Propagating the curvature interval through the mass constraint gives
$\delta\densrest = [f/(1-f)]\times 259 =
\SI{35.5}{\kilo\gram\per\cubic\metre}$ at $f = 0.1205$, against
$\densbulk - \densrest = \SI{69.9}{\kilo\gram\per\cubic\metre}$, so
$\Sstat = 2.0$ for the neck model against $5.8$ for the head model. On
the statistic of Equation~\eqref{eq:sstat}, which is the one that does
not reward a model merely for being stable, the head model is favoured by
a factor of nearly three.

The amplification factor differs sharply between the two models and must
be quoted with either. Because the neck model places less excess mass,
$\ampfac = \densrest/(\densbulk-\densrest) = 1943/69.9 = 27.8$, against
$12.0$ for the head model. A reader comparing the fractional stability of
$\densrest$ between the two models without dividing by $\ampfac$ would
conclude that the neck model was more than twice as well constrained,
when the opposite is the case.

One consequence deserves emphasis rather than burial. The two models
return background densities of $1858$ and
$\SI{1943}{\kilo\gram\per\cubic\metre}$, differing by
$\SI{85}{\kilo\gram\per\cubic\metre}$, or $3.2$ times the combined
uncertainty on the head-model value. The quantity that
Section~\ref{sec:constrains} identifies as the one the method constrains
is therefore not invariant under a change of region shape, only under a
change of plane position within a fixed region shape. That is a narrower
claim than the stability of Table~\ref{tab:itokawasweep} taken alone
would suggest, and we make it in that narrower form.

\subsection{Insensitivity to the assumed mass}
\label{sec:itokawa:mass}

The total mass enters as a constraint, so it is worth asking how much of
the result depends on it. We repeated the reference inversion with
$\Mtot$ varied over a factor of four, from half to twice the measured
value---far beyond the published uncertainty on $GM$, which is about five
per cent \citep{abe2006}. The rotation period was held fixed. Across that
range the recovered contrast moved by $11.7\%$, from $1.404$ to $1.569$,
and the improvement by $0.75$ percentage points, from $16.30\%$ to
$17.05\%$.

Two things should be said about this insensitivity, because it is easy to
over-read. It is not surprising: the contrast is a ratio, and to first
order a change in $\Mtot$ rescales both densities together. What does
\emph{not} rescale is the amplification factor, because
Equation~\eqref{eq:ampsensitivity} makes $\ampfac$ respond to a shift in
$\densbulk$ with a gain of about thirteen. A one per cent error in the
mass therefore moves $\ampfac$ by roughly thirteen per cent, and with it
the division between the algebraic and physical parts of the stability
ratio reported in Section~\ref{sec:constrains}. The contrast is robust to
the assumed mass; the interpretation of the stability is not, and the
error budget of Section~\ref{sec:constrains} carries the mass term for
that reason.

\subsection{Sensitivity to the shape model volume}
\label{sec:itokawa:scale}

Our mesh encloses $\Vtot = \SI{0.0177856}{\cubic\kilo\metre}$ against the
$(1.84\pm0.092)\times10^{7}~\mathrm{m^{3}}$ adopted by \citet{abe2006}, a
difference of $-3.3\%$ and $-0.67\sigma$
(Section~\ref{sec:methods:shapes}). Since $\densbulk = \Mtot/\Vtot$
enters Equation~\eqref{eq:amplification} directly, the difference matters
more than its size suggests.

Rescaling the model linearly to the published volume moves the contrast
by $-0.32\%$ and the excess mass by $-0.92\%$, both under one per cent, while
$\densbulk$ falls from $2012.9$ to
$\SI{1945.7}{\kilo\gram\per\cubic\metre}$ and $\ampfac$ moves from $12.0$
to $12.4$. The rotation period was held at its measured value, so the
rescaling alters the rotational term relative to the gravitational one: it
is a physical change to the problem, not a relabelling of it.
The contrast and the excess mass are robust; $\ampfac$, and hence the
decomposition of Section~\ref{sec:constrains}, is not. We report both
values of $\ampfac$ wherever the decomposition is quoted.

\subsection{Summary of this section}
\label{sec:itokawa:summary}

Itokawa returns a dispersion improvement of $16.9\%$ at the detected
neck, with a head density of
$2800 \pm \SI{129}{\kilo\gram\per\cubic\metre}$ against a background of
$1858 \pm \SI{27}{\kilo\gram\per\cubic\metre}$, a contrast of
$1.51 \pm 0.09$, an excess mass of $7.7 \pm 1.3\%$ of the total, and a
centre-of-mass offset of $\SI{16.9+-2.9}{\metre}$. The background density
is displaced from the bulk density by $5.8$ times its own uncertainty.

The result is stable under decimation at the level the scan grid can
resolve, and the residual resolution sensitivity enters through the plane
detector rather than the objective. It is stable under a change of
dividing plane in the sense that $\densrest$ varies by a quarter of a per
cent across thirteen planes, but not in the sense that $\denshead$ does,
which varies by twenty. It is not stable under a change of region shape:
the neck model returns a background density
$\SI{85}{\kilo\gram\per\cubic\metre}$ higher, though with an improvement
inside the resolution envelope and an $\Sstat$ of $2.0$ against $5.8$, so
the head model is preferred on the paper's own criteria. It agrees with
published determinations, but three of the four rest on the same physics
as ours, and the fourth has intervals wide enough that agreement is weak
evidence. And the objective does not choose its own plane: left to itself
it runs to a head density the meteorite collection forbids, and an
external bound at $\SI{194+-11}{\metre}$ is what stops it.